% SPDX-FileCopyrightText: 2026 Will Cook
% SPDX-License-Identifier: Apache-2.0
%
% A mathematician-facing exposition of the Erdos249257 Lean
% development. Every formalized statement links to the exact Lean source line
% at the pinned commit. Build with `tectonic erdos249-257-exposition.tex` or `make`.
\documentclass[11pt]{article}

\usepackage[T1]{fontenc}
\usepackage[utf8]{inputenc}
\usepackage[a4paper,margin=1in]{geometry}
\usepackage{amsmath,amssymb,amsthm}
\usepackage{booktabs}
\usepackage{seqsplit}
\usepackage{enumitem}
\usepackage[dvipsnames]{xcolor}
\usepackage[colorlinks=true,linkcolor=black,urlcolor=RoyalBlue,citecolor=black]{hyperref}

% ---------------------------------------------------------------------------
% Drill-down machinery: each anchor links to the GitHub blob at a pinned commit.
% ---------------------------------------------------------------------------
\newcommand{\commit}{59f9cc4ce55b2200ca9cbd7ca90b1c0f6894d821}
\newcommand{\repobase}{https://github.com/wcook04/plectis-lean-erdos249-257/blob/\commit}
\newcommand{\PK}{Erdos249257}
% \idn: typeset a Lean identifier (underscores literal, breakable across lines).
\newcommand{\idn}[1]{\texttt{\expandafter\seqsplit\expandafter{\detokenize{#1}}}}
% The source coordinates remain in every hyperlink target and in the verifier's
% manifest.  The paper shows a declaration name for principal results and a
% quiet source marker for secondary links, so provenance does not dominate prose.
\newcommand{\lref}[3]{\href{\repobase/\PK/#1\#L#2}{\idn{#3}}}
\newcommand{\lrefx}[3]{\href{\repobase/\PK/#1\#L#2}{\texttt{Lean source}}}
\newcommand{\lloc}[2]{\href{\repobase/\PK/#1\#L#2}{\texttt{Lean source}}}
% \rref{rootfile}: link a repository-root file (underscores handled).
\newcommand{\rref}[1]{\href{https://github.com/wcook04/plectis-lean-erdos249-257/blob/main/\detokenize{#1}}{\texttt{\detokenize{#1}}}}
\emergencystretch=2em

\newtheorem{result}{Result}
\newtheorem{proposition}[result]{Proposition}
\theoremstyle{definition}
\newtheorem{definition}[result]{Definition}
\theoremstyle{remark}
\newtheorem*{remark}{Remark}

\newcommand{\N}{\mathbb{N}}
\newcommand{\Z}{\mathbb{Z}}
\newcommand{\Q}{\mathbb{Q}}
\newcommand{\R}{\mathbb{R}}
\newcommand{\ph}{\varphi}
\newcommand{\lcm}{\operatorname{lcm}}
\newcommand{\Irr}{\text{ irrational}}
\newcommand{\Qzero}{Q_0}

\title{\vspace{-2em}Erd\H{o}s Problems 249 and 257 in Lean~4:\\
{\large formalized results, certificate reductions, and open questions}}
\author{Will Cook}
\date{July 2026}

\begin{document}
\sloppy
\maketitle

\begin{abstract}
\noindent
We formalize in Lean~4 several results around Erd\H{o}s Problems 249 and 257.
For every integer base $b\ge2$, the development verifies the irrationality of
the Erd\H{o}s--Borwein series $\sum_{n\ge1}(b^n-1)^{-1}$, together with several
structured infinite-support variants. For the totient constant
$S=\sum_{n\ge1}\ph(n)/2^n$, whose irrationality remains open, we prove that any
rational representation has denominator greater than
$\Qzero\approx7.96\times10^{34}$, and formalize an exact reduction of
irrationality to an unbounded family of finite decidable certificates. We also
describe the proof architecture and auxiliary rationality criteria based on
binary carry systems. This work solves neither Erd\H{o}s~\#249 nor the universal
form of Erd\H{o}s~\#257.
\end{abstract}

\section{Introduction}\label{sec:intro}

Two open questions provide the setting. Erd\H{o}s~\#249 asks whether the
totient constant
\[
S=\sum_{n\ge1}\frac{\ph(n)}{2^n}
\]
is irrational. Erd\H{o}s~\#257 asks whether
$\sum_{n\in A}(2^n-1)^{-1}$ is irrational for every infinite
$A\subseteq\N$. The second problem contains the classical full-support
Erd\H{o}s--Borwein constant, but the passage from structured supports to an
arbitrary infinite support remains open.

The development has four principal outputs. First, it formalizes the
irrationality of $\sum_{n\ge1}(b^n-1)^{-1}$ for every integer $b\ge2$, and of
several structured infinite-support variants. Second, it places $S$ on the
same Mersenne--Lambert ladder and proves useful positive, signed, and
probabilistic representations. Third, it gives the unconditional denominator
bound $q>\Qzero$ for any rational expression $S=p/q$. Fourth, it proves an
exact finite-certificate reduction: irrationality follows if the certified
non-integrality witnesses occur at unbounded parameters, and the existence of
a witness at fixed parameters is equivalent to the corresponding tail
difference being non-integral.

\begin{center}
\fbox{\begin{minipage}{0.91\textwidth}
\textbf{Scope.} The irrationality of $S$ and the universal infinite-support
statement remain open. The denominator theorem and certificate reduction are
genuine unconditional and conditional advances toward \#249, respectively;
neither is presented as a solution. The support theorems establish named
families toward \#257, not the universal quantifier.
\end{minipage}}
\end{center}

The formalization-specific contribution is not merely a list of declarations.
It isolates a common architecture in which an infinite irrationality question
is converted into finite arithmetic checked by the Lean kernel, while the
remaining analytic uniformity statement is exposed rather than hidden inside
computation. Exact links from the informal statements to the pinned source are
kept as quiet hyperlinks; the complete declaration map is in
Appendix~\ref{app:index}.

The prior-art boundary is theorem-specific. The full-support theorem is due to
Erd\H{o}s~\cite{erdos1948}, and the pairwise-coprime support condition is also
classical~\cite{erdos1968}. Periodic-coefficient Lambert series have a separate
literature~\cite{lucatachiya}. Kova\v{c}--Tao record the strict-tail geometry
used by the achievement-set appendix~\cite{kovactao}, while Wang--Grau Ribas
use the integral carry recurrence in the weighted binary setting
\cite{wanggrauribas}. We make no priority claim for these mathematical
ingredients; the contribution here is their verified integration, the stated
certificate reductions and bounds, and the formalization lessons described
below.

\begin{center}
\fbox{\begin{minipage}{0.94\textwidth}
\centering
\textbf{The Mersenne--Lambert ladder.}\quad $\displaystyle L(f)=\sum_{n\ge1}\frac{f(n)}{2^n-1}$\\[6pt]
\begin{tabular}{@{}cccc@{}}
$L(\mu)=\tfrac12$ & $L(\ph)=2$ & $L(\mathbf 1)=E$ & $L(\ph*\mu)=S$\ (\#249)\\
rational & rational & \textbf{irrational (proved)} & \textbf{open}\\
\end{tabular}\\[10pt]
\textbf{Logical status of the \#249 reduction.}\\[3pt]
$S\in\Q\ \Longrightarrow\ \text{eventual binary periodicity}
\ \Longrightarrow\ R_{N+h}-R_N\in\Z$ for the eventual period.\\[3pt]
$\textsf{certifiedKill}(h,N,L)\ \Longrightarrow\ R_{N+h}-R_N\notin\Z$, and\\[-2pt]
$\exists L\,\textsf{certifiedKill}(h,N,L)\ \Longleftrightarrow\ R_{N+h}-R_N\notin\Z$.\\[4pt]
{\footnotesize Solid implications and the equivalence are proved; the
unresolved step is an unbounded supply of such non-integral tail differences.}
\end{minipage}}
\end{center}

\section{The Mersenne--Lambert ladder}\label{sec:ladder}

Define the Lambert-type transform of an arithmetic function $f$ by
\[
L(f)\;=\;\sum_{n\ge1}\frac{f(n)}{2^n-1}\;=\;\sum_{m\ge1}\frac{(f*\mathbf 1)(m)}{2^m},
\qquad (f*\mathbf 1)(m)=\sum_{d\mid m} f(d),
\]
the second equality being the standard Lambert rearrangement (one Dirichlet
convolution by the constant $\mathbf 1=\zeta$). Walking the ladder by convolution
gives five values with rational, irrational, open, or transcendental status.

\begin{center}
\begin{tabular}{@{}llll@{}}
\toprule
input $f$ & $L(f)$ & value & status in the repository\\
\midrule
$\mu$ (M\"obius) & $L(\mu)$ & $1/2$ & rational, proved\\
$\ph$ (totient) & $L(\ph)$ & $2$ & rational, proved\\
$\mathbf 1$ (constant one) & $L(\mathbf 1)$ & $E$ & \textbf{irrational, machine-checked}\\
$A=\ph*\mu$ & $L(A)$ & $S$ & \textbf{open (Erd\H{o}s \#249)}\\
$\mathrm{Id}$ & $L(\mathrm{Id})$ & $\sum\sigma(m)/2^m$ & transcendental (Nesterenko~1996~\cite{nesterenko}, cited only)\\
\bottomrule
\end{tabular}
\end{center}

The weight $A=\ph*\mu$ is the primitive-conductor count: $A(p)=p-2$ on primes,
$A\ge0$, and $A*\zeta=\ph$. One convolution by $\zeta$ separates the open
constant $S$ from the rational $2$, and one more separates rational from
transcendental ($\ph*\zeta=\mathrm{Id}$, $\mathrm{Id}*\zeta=\sigma$). So $S$ is
the rung one convolution from the machine-checked irrational $E$: the same series
$\sum(\cdot)/(2^d-1)$, with the constant weight $1$ replaced by the nonnegative
weight $A$.

\begin{result}[the rational rungs]
$\sum_{d\ge1}\mu(d)/(2^d-1)=1/2$ and $\sum_{d\ge1}\ph(d)/(2^d-1)=2$, exactly.
\end{result}
\noindent Formalized: \lref{MersenneLambertLadder.lean}{527}{tsum_moebius_div_two_pow_sub_one_eq_half}
and \lref{MersenneLambertLadder.lean}{515}{tsum_totient_div_two_pow_sub_one_eq_two}; restated on the
\#249 series shape at \lrefx{CertificateKernel.lean}{18105}{tsum_moebius_div_two_pow_sub_one_eq_half}
and \lrefx{CertificateKernel.lean}{18100}{tsum_totient_div_two_pow_sub_one_eq_two}.

\begin{result}[the positive lift, $L(A)=S$]\label{res:lift}
With the primitive-conductor weight $A=\ph*\mu$,
\[
\sum_{d\ge1}\frac{A(d)}{2^d-1}\;=\;\sum_{n\ge1}\frac{\ph(n)}{2^n}\;=\;S.
\]
\end{result}
\noindent Formalized: \lref{CertificateKernel.lean}{18116}{tsum_primWeight_div_two_pow_sub_one_eq_totient_series},
with the weight $A$ defined at \lrefx{MersenneLambertLadder.lean}{247}{primWeight}. This
places $S$ in the same $\sum(\cdot)/(2^d-1)$ family as $E$, with a nonnegative weight.

\begin{result}[the M\"obius-square lens]\label{res:lens}
\[
S \;=\; \tfrac12 \;+\; \sum_{d\ge1}\frac{\mu(d)}{(2^d-1)^2}.
\]
\end{result}
\noindent Formalized: \lref{CertificateKernel.lean}{18125}{totient_series_eq_half_add_moebius_mersenne_square}.
Writing $T=\sum_{d\ge1}\mu(d)/(2^d-1)^2$, irrationality of $S$ is equivalent to
irrationality of $T$, since $S=\tfrac12+T$. The factor $1/(2^d-1)^2$ is the
probability that $d$ divides two independent fair-coin waiting times
(Section~\ref{sec:geom}). The signed constant, the coprime-pair mass, and the
tail differences of Section~\ref{sec:reduction} are three readings of the same
quantity.

\section{Erd\H{o}s--Borwein-type irrationality (the \#257 direction)}\label{sec:eb}

\subsection{Full support: the Erd\H{o}s--Borwein constant, every base}

\begin{result}[full-support irrationality]\label{res:full}
For every integer $b\ge2$, the series $\displaystyle\sum_{n\ge1}\frac{1}{b^n-1}$ is
irrational. In particular the Erd\H{o}s--Borwein constant $E=\sum_{n\ge1}1/(2^n-1)$
is irrational.
\end{result}
\noindent Formalized: \lref{CertificateKernel.lean}{7999}{irrational_erdosSum_full_support}
(all bases) and \lref{CertificateKernel.lean}{8006}{irrational_erdosBorwein_series}
(the base-$2$ corollary). This is Erd\H{o}s's 1948 theorem~\cite{erdos1948}. In the ladder of
Section~\ref{sec:ladder} it is the rung $L(\mathbf 1)$, and it is the fixed
machine-checked point one convolution away from the open $S$.

The proof runs through the certificate machine of the kernel rather than through a
direct digit argument: a bounded Bertrand/CRT frame on a first block, a
middle-window average over divisor pairs with a pigeonhole selection, and an
explicit closure of the parameters. The Lambert bridge
$\sum 1/(b^n-1)=\sum_m \tau(m)/b^m$ (divisor-count series) is
\lrefx{CertificateKernel.lean}{5864}{erdosSum_full_support_eq_tsum_divisor_count} in the source.

\subsection{Named infinite-support cases}

Beyond full support, the development proves irrationality for a family of infinite
supports $A$, at every base $b\ge2$. Each is a formalized case of the \#257
statement family. None is the universal statement, and none is claimed as new
mathematics.

\begin{result}[support-class family]\label{res:support}
For every integer $b\ge2$, the series $\sum_{n\in A} 1/(b^n-1)$ is irrational for
each of the following supports $A$:
\begin{enumerate}[label=\textup{(\alph*)},leftmargin=2em]
\item factorial support $A=\{n! : n\ge1\}$, giving $\sum 1/(b^{n!}-1)$;
\item power-of-two support $A=\{2^n : n\ge0\}$, giving $\sum 1/(b^{2^n}-1)$;
\item multiples $A=d\N$;
\item pairwise-coprime supports with summable reciprocals (Erd\H{o}s's
condition~\cite{erdos1968});
\item eventually-periodic supports, residue classes, and the odd numbers.
\end{enumerate}
\end{result}
\noindent These sort into three groups. The classical case is (d): infinite
pairwise-coprime $A$ with $\sum_{a\in A}1/a<\infty$, formalized from Erd\H{o}s's
condition. The lcm-gap named cases are (a) and (b), whose support gaps outgrow the
running $\lcm$; the structured families are (c) and (e).\footnote{Periodic-coefficient
Lambert series have their own literature; see e.g.\ Luca--Tachiya~\cite{lucatachiya}.
The cases here are formalized in the repository's own framework and are not claimed
as new results.} Formalized: (a)
\lrefx{CertificateKernel.lean}{5706}{irrational_erdosSum_factorial_support};
(b) \lrefx{CertificateKernel.lean}{5730}{irrational_erdosSum_two_pow_support};
(c) \lrefx{CertificateKernel.lean}{8774}{irrational_erdosSupportSeries_multiples};
(d) \lrefx{CertificateKernel.lean}{10447}{irrational_erdosSupportSeries_pairwise_coprime};
(e) \lrefx{CertificateKernel.lean}{11275}{irrational_erdosSupportSeries_eventuallyPeriodic},
\lrefx{CertificateKernel.lean}{11343}{irrational_erdosSupportSeries_residueClass},
\lrefx{CertificateKernel.lean}{11357}{irrational_erdosSupportSeries_odd}. The two
base-$2$ instances that name the problem directly are
\lrefx{CertificateKernel.lean}{5753}{erdos257_family_factorial_instance} and
\lrefx{CertificateKernel.lean}{5761}{erdos257_family_two_pow_instance}.

The factorial and power-of-two cases are worth a word because they are not
reachable by a Liouville shortcut: the power-of-two gap $2^k-\lcm$ stays within a
bounded ratio of the $\lcm$, so the terms do not decay fast enough for a
transcendence-measure argument, and a uniform criterion is needed. The mechanism
is what the repository calls \emph{period noncollapse}: a
prime-valuation-deficit witness forces the multiplicative order of the base modulo
the relevant modulus not to collapse, which rules out the long repeated digit
blocks a rational value would require.

\begin{remark}
A source comment near \lrefx{CertificateKernel.lean}{5943}{lcm_gap_hypothesis_fails_full_support}
describes a second, near-integer route to full support whose witnesses are not
built there, and says so. The full-support theorem is proved by the separate
certificate engine cited under Result~\ref{res:full}. The two coexist; the
near-integer route is one proof strategy, not a retraction of the theorem.
\end{remark}

\subsection{What stays open}

The universal Erd\H{o}s \#257 statement, irrationality of $\sum_{n\in A}1/(2^n-1)$
for \emph{every} infinite $A$, is not proved. The proved cases are a family, not
the quantified theorem, and the repository marks this in code
(the non-claim \idn{not_erdos_257_solution} in \rref{NON_CLAIMS.md}).

\section{The totient constant $S$ (Erd\H{o}s \#249)}\label{sec:249}

Whether $S=\sum_{n\ge1}\ph(n)/2^n$ is irrational is open. The development gives two
things and claims neither as a solution: an unconditional exclusion of
small-denominator rationals, and a conditional reduction to finite certificates.

\subsection{Unconditional: no small denominator}\label{sec:uncond}

\begin{result}[denominator exclusion]\label{res:farey}
Set
\[
\Qzero:=79\,639\,646\,646\,701\,375\,323\,355\,774\,875\,831\,053
       \approx7.96\times10^{34}.
\]
For every integer $p$ and every $q$ with $1\le q\le\Qzero$,
\[
S \;\ne\; \frac{p}{q}.
\]
Equivalently, if $S$ is rational then its reduced denominator exceeds that bound.
\end{result}
\noindent Formalized: \lloc{CertificateKernel.lean}{18055} (\,\idn{tsum_totient_div_pow_two_ne_ratCast_of_den_le_79639646646701375323355774875831053}\,).
The bound is reached by a ladder of finite, elementary exclusions,
\[
4838 \;\to\; 4\,194\,304 = 2^{22} \;\to\; 2.49\times10^{17}\ (\text{Farey } K{=}120)
\;\to\; 7.96\times10^{34}\ (\text{Farey } K{=}240),
\]
each a mediant/Farey argument that traps any candidate $m/q$ inside a forbidden
gap. See \lrefx{CertificateKernel.lean}{15346}{tsum_totient_div_pow_two_ne_ratCast_of_den_le_4838}
for the first rung and \lrefx{GapFareyBound.lean}{80}{gap_check_window_1_120_le_248672326362367909},
\lrefx{GapFareyBound.lean}{156}{gap_check_window_1_240_le_79639646646701375323355774875831053}
for the two Farey windows.

This is an explicit lower bound on the denominator of any rational
representation, not a proof of irrationality. A rational with a larger
denominator is not excluded by it.

\begin{remark}[one bound, four representations]
Because the ladder identities of Section~\ref{sec:ladder} are equalities, the same
finite record transfers verbatim to three other representations of the constant:
to the positive lift $\sum_d A(d)/(2^d-1)$
(\lloc{CertificateKernel.lean}{18144}),
to the signed M\"obius-square constant $T=\sum_d\mu(d)/(2^d-1)^2$ at half the bound
(\lloc{CertificateKernel.lean}{18158}),
and to the coprimality probability of Section~\ref{sec:geom} at half the bound.
Half appears because $S=\tfrac12+T$ turns a denominator $d$ for $T$ into $2d$ for $S$.
\end{remark}

\subsection{The geometric picture: $S$ as coprime-pair mass}\label{sec:geom}

Let $X,Y$ be independent fair-coin waiting times, $\Pr(X=n)=2^{-n}$ for $n\ge1$.
Then $\Pr(d\mid X)=1/(2^d-1)$, so the Lambert transform is a divisor calculus of
$X$: $L(f)=\mathbb E[(f*\zeta)(X)]$. The rungs read as $L(\mu)=\Pr(X=1)=1/2$,
$L(\ph)=\mathbb E[X]=2$, $L(\mathbf 1)=\mathbb E[\tau(X)]=E$, and $L(A)=\mathbb E[\ph(X)]=S$.

Counting visible lattice points $(a,b)$ with $a+b=n$, $a>0$, $\gcd(a,b)=1$ gives
exactly $\ph(n)$, uniformly in $n$, so $S$ is itself a coprime-pair mass.

\begin{result}[coprimality probability]\label{res:coprime}
\[
S \;=\; \tfrac12 \;+\; \Pr\!\big(\gcd(X,Y)=1\big)
\;=\; \tfrac12 \;+\!\!\sum_{\substack{a,b\ge1\\ \gcd(a,b)=1}}\!\! 2^{-(a+b)}.
\]
\end{result}
\noindent Formalized: \lref{CertificateKernel.lean}{18228}{totient_series_eq_half_add_visible_coprime_pairs},
with the underlying lattice identity at
\lrefx{CertificateKernel.lean}{18215}{tsum_visible_coprime_pairs_eq_totient_series}.
Base $2$ is the one point where $\Pr(X=a)\Pr(Y=b)=2^{-(a+b)}$ needs no normalizing
constant, so the totient sum and the coprimality probability line up with no
boundary correction. In this reading, Erd\H{o}s \#249 asks whether two independent
fair-coin waiting times are coprime with irrational probability.

\subsection{Conditional: reduction to finite certificates}\label{sec:reduction}

The reduction turns irrationality of $S$ into the non-vanishing of a decidable
predicate at infinitely many parameters. Everything in it is elementary:
$\ph(n)\le n$, geometric tails, and integer window arithmetic. There is no q-Pad\'e
input, no Stern--Brocot input, and no unformalized citation
(\lrefx{TotientTailPeriodKiller.lean}{47}{TotientTailPeriodKiller}).

\paragraph{Step 1: rationality forces an integer tail-period law.} Write the
fractional layer of $2^N S$ as the tail
\[
R_N \;=\; \sum_{m\ge1}\frac{\ph(N+m)}{2^m}.
\]
If $S$ is rational, its binary expansion is eventually periodic, which forces, for
a suitable period $h$ and all large $N$, the integrality
$R_{N+h}-R_N\in\Z$. To refute integrality at a finite depth $L$, form the window
discrepancy
\[
\Delta_{h,N,L}=\sum_{j=0}^{L-1}\big(\ph(N+h+1+j)-\ph(N+1+j)\big)\,2^{\,L-1-j}\in\Z,
\]
the first $L$ binary digits of $R_{N+h}-R_N$ up to a controlled deep tail (bounded
using only $\ph(n)\le n$), and call the pair \emph{certified} when the residue of
$\Delta_{h,N,L}$ modulo $2^L$ misses a neighbourhood of $0$:
\[
\textsf{certifiedKill}(h,N,L):\quad
N+h+L+2 \;<\; \big(\Delta_{h,N,L}\bmod 2^L\big) \;<\; 2^L-(N+h+L+2).
\]
This is a finite integer computation, decidable by the kernel. It is the natural
finite object here because a rational $S$ would pin the tail difference to an
integer, and a residue that stays a fixed distance from $0$ is a finite proof that
it did not.

\begin{result}[tail-period reduction]\label{res:red21}
If for every period $h\ge1$ and every threshold $N_0$ there exist $N\ge N_0$ and
$L$ with $\textsf{certifiedKill}(h,N,L)$, then $S$ is irrational.
\end{result}
\noindent Formalized: \lref{TotientTailPeriodKiller.lean}{386}{irrational_totient_series_of_certificate_supply};
the tail $R_N$ and the predicate are
\lrefx{TotientTailPeriodKiller.lean}{57}{totientTail} and
\lrefx{TotientTailPeriodKiller.lean}{72}{certifiedKill}.

\paragraph{Step 2: collapse the two parameters to one.} Every multiple of a period
is a period, so $\lcm(1,\dots,t)$ is the universal-period envelope: any eventual
period of a rational $S$ divides it once $t$ is large. Standing on that one
diagonal ray covers all possible periods at once, and the position parameter
drops out too.

\begin{result}[diagonal reduction]\label{res:diag}
Write $\lcm(1..t)$ for $\lcm(1,\dots,t)$. If for every $t_0$ there exist $t\ge t_0$
and $L$ with $\textsf{certifiedKill}\big(\lcm(1..t),\,\lcm(1..t),\,L\big)$, then $S$
is irrational.
\end{result}
\noindent Formalized: \lref{LcmDiagonalReduction.lean}{92}{irrational_totient_series_of_lcm_diagonal_certificate_supply}.

\paragraph{Step 3: the cone, flatness, and completeness.} On the cone
$\{k\cdot\lcm(1..t)\}$, rationality forces not one pairwise relation but a single
fractional constant across the whole cone.

\begin{result}[cone flatness]\label{res:flat}
If $S$ is rational, then the relevant tail values share one fractional part across
the whole $\lcm$ cone.
\end{result}
\noindent Formalized: \lref{LcmConeFlatness.lean}{90}{rational_totient_series_forces_lcm_cone_flatness}.

The certificates are exact, not a lossy heuristic: a certificate exists exactly
when the object it certifies is non-integral. Thus the finite computations are
machine-checked witnesses rather than numerical experiments.

\begin{result}[certificate completeness]\label{res:complete}
For all $h,N$,
\[
\big(\exists L,\ \textsf{certifiedKill}(h,N,L)\big)
\quad\Longleftrightarrow\quad
R_{N+h}-R_N\notin\Z.
\]
\end{result}
\noindent Formalized: \lref{LcmConeFlatness.lean}{316}{exists_certifiedKill_iff_tail_diff_notMem_int}.
Finite certificates witness non-integrality. A joint refuter
can interrogate a menu of cone vertices at once and is sharper than any pairwise
test: it refutes the single-fractional-part model that rationality predicts on a
finite sample of cone vertices
(\lref{LcmConeNonflat.lean}{126}{exists_nonintegral_pair_of_coneNonflatCert}).

\paragraph{The certificates are non-empty.} The initial segment of the required
supply is machine-checked, so the reduction is not vacuous. By completeness, each
of these is a kernel-checked proof of a specific non-integral tail difference,
not a numerical experiment.

\begin{result}[verified finite instances]\label{res:deposits}
The following are checked by the Lean kernel:
\begin{itemize}[leftmargin=1.4em]
\item $\textsf{certifiedKill}(h,12,16)$ for every period $h\in[1,8]$;
\item certified kills for every period $h\in[1,16]$ at a fixed window;
\item diagonal certificates $\textsf{certifiedKill}(\lcm(1..t),\lcm(1..t),L)$ for $t\in[1,6]$, and at $t=7,8$;
\item a tabulated joint cone-vertex certificate.
\end{itemize}
\end{result}
\noindent Formalized: \lrefx{TotientTailPeriodKiller.lean}{396}{certifiedKill_all_small},
\lrefx{CarrySurvivorExtinction.lean}{324}{certifiedKill_all_upto_sixteen},
\lrefx{LcmDiagonalReduction.lean}{217}{certifiedKill_periodLcm_diagonal_upto_six}
(with \lrefx{LcmDiagonalReduction.lean}{222}{certifiedKill_periodLcm_diagonal_at_seven},
\lrefx{LcmDiagonalReduction.lean}{226}{certifiedKill_periodLcm_diagonal_at_eight}),
and \lrefx{LcmConeNonflat.lean}{506}{coneNonflatCert_cells}. These instances show
that each layer of the reduction is non-vacuous and, by completeness, exact. They
are not evidence that the unbounded supply exists.

\paragraph{What is open.} The reduction converts \#249 into the statement that the
certificate supply is \emph{unbounded}: certificates exist for arbitrarily large
parameters. By completeness (Result~\ref{res:complete}) this is equivalent to the
tail differences being non-integral infinitely often, and equivalently to the
Farey gap denominators of Section~\ref{sec:uncond} being unbounded: the same
obstruction seen as a rational denominator, an eventual binary period, and a flat
tail on the $\lcm$ cone. That unboundedness is the remaining analytic question,
and it is left open. The reduction is real; the closure is not asserted.

\section{Formalization architecture and mathematical lessons}\label{sec:architecture}

The formal development repeatedly separates a finite, kernel-checkable core
from the analytic statement that would have to hold uniformly. For the
full-support Erd\H{o}s--Borwein theorem, the core is a bounded certificate
engine assembled from a Bertrand/CRT frame, divisor-pair averaging, and an
explicit parameter closure. For \#249, the same design principle appears more
sharply: eventual binary periodicity predicts an integral tail difference;
finite modular arithmetic can refute that prediction at any fixed pair of
parameters; and completeness proves that this computation loses no
information. Lean therefore makes the open boundary visible: the missing
ingredient is not another finite verification but a theorem producing
witnesses at unbounded parameters.

Three changes of coordinates serve different mathematical purposes. The
Mersenne--Lambert transform places $S$ next to the Erd\H{o}s--Borwein constant
and exposes the nonnegative primitive-conductor weight. The M\"obius-square
identity turns the same number into a signed rapidly convergent series. The
coprime-pair identity turns it into a probability. Formalizing all three and
proving that the denominator exclusion transfers between them prevents an
informal change of representation from silently changing the claim.

The auxiliary carry modules ask what rationality would force rather than
claiming a new irrationality theorem. They give an exact tempered-orbit
criterion, Boolean--M\"obius coordinates for support series, constraints from
residue wraps and reciprocal mass, and a greedy description of the value set.
Together they rule out several tempting but false proof strategies---for
example, rationality need not give a bounded carry alphabet---and identify the
kind of infinite Boolean carry path that a future proof would have to exclude.
The precise statements and source links are retained in
Appendix~\ref{sec:carry}; their supporting role is why they are not a second
main narrative here.

This architecture is reusable beyond the two named problems: choose a
representation in which rationality imposes a rigid tail law, isolate a finite
predicate that soundly refutes that law, prove the predicate complete at fixed
parameters, and state the remaining uniformity theorem separately. The last
step matters as much as the first three. It prevents a large verified finite
range from being mistaken for the infinite theorem.

\section{Artefact availability and verification}\label{sec:verify}

The Lean~4 source, exact informal-to-formal links, and build metadata are
available in the public repository
\href{https://github.com/wcook04/plectis-lean-erdos249-257}{\texttt{plectis-lean-erdos249-257}}.
This paper is linked to commit \href{\repobase}{\texttt{59f9cc4}}; each blue
Lean link targets an exact declaration at that commit (for example, the
base-$2$ corollary in Result~\ref{res:full} resolves to
\lrefx{CertificateKernel.lean}{8006}{irrational_erdosBorwein_series}). The package uses
\texttt{leanprover/lean4:v4.29.1} and the Mathlib revision pinned by
\href{\repobase/lake-manifest.json}{\texttt{lake-manifest.json}}.

There is no \texttt{sorry} or \texttt{admit} in the package. Finite
computations use \texttt{decide}, whose proof terms are checked by the Lean
kernel; the package does not use \texttt{native\_decide}. Reproduction requires
only the pinned toolchain and the ordinary project build:
\begin{verbatim}
lake exe cache get
lake build
\end{verbatim}
The second command elaborates every declaration and kernel-checks every proof
term against the pinned Mathlib dependency. Continuous integration executes
the same build. Appendix~\ref{app:index} gives a compact map of the principal
informal statements to declarations; the repository retains the exhaustive
machine-readable provenance checked by the paper-link verifier.

\section{Conclusion}\label{sec:conclusion}

The formalization fixes two mathematical landmarks and two exact frontiers. It
verifies the all-integer-base Erd\H{o}s--Borwein theorem and several structured
support cases toward \#257. For \#249 it proves that a rational value of $S$
would have denominator exceeding $\Qzero$, and it reduces irrationality to an
unbounded supply of finite certificates whose fixed-parameter semantics are
complete. The unresolved \#249 step is precisely that unboundedness; the
unresolved \#257 step is the passage from the formalized support classes to
every infinite support. Neither frontier is softened by the size of the
verified artefact.

The main lesson from formalization is the alignment of representations and
proof obligations. The Lambert, M\"obius, probabilistic, periodic-tail, and
carry descriptions are not competing stories: each makes a different
rationality obstruction explicit. A future advance must exploit one of those
obstructions uniformly, not merely extend the finite range already checked.

\appendix
\section{Auxiliary rationality criteria from binary carry systems}\label{sec:carry}

Four further modules provide additional criteria. They prove no new
irrationality. Instead they pin down, in machine-checked form, what a rational
value of the series in question would have to look like: as an integer carry
orbit, as a M\"obius-inverted support coordinate, as a residue-wrap cycle
modulo the odd part of the denominator, and as a point of a greedy subset-sum
geometry. Everything in this section is an equivalence or an unconditional
constraint on the rational side; nothing excludes a rational value of either
open series, and each module states that boundary in its header. The prior-art
boundary is theorem-family-specific. The integral carry recurrence has a
direct antecedent in Wang--Grau Ribas~\cite{wanggrauribas}; the strict-tail
geometry is recorded by Kova\v{c}--Tao~\cite{kovactao}; M\"obius inversion,
repetend algebra, and divisor averaging are classical. We make no claim of
mathematical priority for the converse/rigidity, certificate-normal-form, or
coupled reciprocal-mass/collision statements collected here.

\subsection{Tail-orbit rigidity: a rationality criterion}

For a coefficient sequence $c:\N\to\N$ with $c(n)\le n$, write
\[
X_c=\sum_{n\ge1}\frac{c(n)}{2^n},
\qquad
T_c(N)=\sum_{j\ge1}\frac{c(N+j)}{2^j},
\]
the series and its scaled tail after $N$ binary digits. Call an integer
sequence $u$ a \emph{tempered orbit} for $c$ with multiplier $v\ge1$ if
\[
u(N+1)=2\,u(N)-v\,c(N+1)\ \text{ for all } N,
\qquad\text{and}\qquad
u(N)/2^N\to0 .
\]
The recurrence is exact binary long division against the coefficient stream;
the boundary condition forbids the homogeneous solution. Wang and Grau Ribas
use the integral carry recurrence forced by rationality for the weighted binary
special case $c(n)=n d_n$~\cite{wanggrauribas}. The module abstracts that
forward mechanism to arbitrary nonnegative $c(n)\le n$; its explicit converse
and subexponential-orbit rigidity are the local additions, with no priority
claim attached.

\begin{result}[tail-orbit rigidity]\label{res:rigidity}
Let $c(n)\le n$ for all $n$. Every tempered orbit is the scaled tail,
$u(N)=v\,T_c(N)$ for all $N$; and $X_c$ is rational if and only if a tempered
orbit exists for some multiplier $v\ge1$.
\end{result}
\noindent Formalized:
\lref{GenericTailOrbitRigidity.lean}{180}{temperedBinaryOrbit_eq_scaledTail}
(rigidity) and
\lref{GenericTailOrbitRigidity.lean}{260}{binaryCoeffSeries_rational_iff_exists_temperedBinaryOrbit}
(the criterion), restated in Mathlib's irrationality vocabulary at
\lrefx{GenericTailOrbitRigidity.lean}{269}{not_irrational_binaryCoeffSeries_iff_exists_temperedBinaryOrbit}.
The engine is one observation: a real orbit with $d(N+1)=2\,d(N)$ and
$d(N)=o(2^N)$ vanishes identically
(\lrefx{GenericTailOrbitRigidity.lean}{159}{doublingOrbit_eq_zero_of_tempered}).

The tempered boundary is essential: the
homogeneous parasite $N\mapsto2^N$ can be added to any orbit without
disturbing the recurrence, so positivity of an orbit certifies nothing by
itself, and positivity is nowhere advertised as an equivalent criterion. The
linear-growth hypothesis covers both constants introduced in
Section~\ref{sec:intro}:
$\ph(n)\le n$ for the \#249 coefficients, and $f_A(n)\le\tau(n)\le n$ for the
\#257 support coefficients $f_A(n)=\#\{a\in A:a\mid n\}$. The modules below
instantiate the criterion on supports; the totient instantiation is not taken
here.

\subsection{Boolean--M\"obius carry coordinates (the \#257 direction)}

At base $2$ the support series of Section~\ref{sec:eb} is itself a binary
coefficient series, $\sum_{n\in A}1/(2^n-1)=X_{f_A}$
(\lrefx{BooleanMobiusCarry.lean}{379}{erdosSupportSeries_two_eq_binaryCoeffSeries}).
Two exact coordinate changes then turn a hypothetical rational value into a
support-free object. The divisor transform and M\"obius inversion below are
classical Lambert-series infrastructure; the contribution here is their
Lean-checked composition with the tempered carry interface.

\begin{result}[Boolean M\"obius inversion, both directions]\label{res:boolmob}
On positive integers, $f_A=\mathbf 1_A*\zeta$ and $\mu*f_A=\mathbf 1_A$, where
$\mathbf 1_A$ is the indicator of $A$. Conversely, every integer-valued
arithmetic function $f$ whose M\"obius transform is Boolean,
$(\mu*f)(n)\in\{0,1\}$ for all $n\ge1$, satisfies $f=f_B$ for the support
$B=\{n\ge1:(\mu*f)(n)=1\}$ selected by that transform.
\end{result}
\noindent Formalized:
\lref{BooleanMobiusCarry.lean}{95}{positiveSupportBitAF_mul_zeta},
\lref{BooleanMobiusCarry.lean}{113}{moebius_mul_supportCoeffAF}, and, with no
search or finiteness hypothesis, the converse
\lref{BooleanMobiusCarry.lean}{190}{eq_supportCoeffAF_booleanMobiusSupport_of_boolean};
the selected support is
\lrefx{BooleanMobiusCarry.lean}{158}{booleanMobiusSupport}.

On the carry side, Result~\ref{res:rigidity} applies at $c=f_A$: the support
series is rational exactly when a tempered carry orbit for $f_A$ exists
(\lrefx{BooleanMobiusCarry.lean}{386}{erdosSupportSeries_rational_iff_exists_temperedCarry});
a displayed fraction $p/q$ corresponds to an orbit $U$ with multiplier $q$ and
initial value $U(0)=p$
(\lrefx{BooleanMobiusCarry.lean}{465}{support_fraction_iff_exists_temperedCarry});
and exact division recovers the coefficient from the orbit,
$\bigl(2U(N)-U(N+1)\bigr)/q=f_A(N+1)$
(\lrefx{BooleanMobiusCarry.lean}{494}{temperedCarry_quotient_eq_supportCoeff}).
The elementary divisor-pair bound $\tau(n)\le2\lfloor\sqrt n\rfloor$
(\lrefx{BooleanMobiusCarry.lean}{218}{card_divisors_le_two_mul_sqrt}) confines
the tails to a square-root strip, $T_{f_A}(N)\le2\sqrt N+4$
(\lrefx{BooleanMobiusCarry.lean}{292}{binaryCoeffTail_supportCoeff_le_two_sqrt_add_four}),
hence every orbit state to $0<U(N)\le q\,(2\sqrt N+4)$
(\lrefx{BooleanMobiusCarry.lean}{504}{temperedCarry_pos_of_exists_pos_mem},
\lrefx{BooleanMobiusCarry.lean}{518}{temperedCarry_le_denominator_mul_two_sqrt_add_four});
and the strip is strong enough to re-derive the tempered boundary, so the
analytic side condition can be traded for one inequality. The two coordinate
changes compose.

\begin{result}[certificate-level equivalence]\label{res:carrycert}
Fix an integer $p$ and an integer $q\ge1$. There is a nonempty support $A$ of
positive integers with $\sum_{n\in A}1/(2^n-1)=p/q$ if and only if there is an
integer sequence $U$ with $U(0)=p$ such that: every $U(N)$ is positive;
$U(N)\le q\,(2\sqrt N+4)$; $q$ divides $2U(N)-U(N+1)$ for every $N$; and the
M\"obius transform of the quotient sequence $\bigl(2U(N)-U(N+1)\bigr)/q$ is
Boolean. The support is not guessed: it is reconstructed from the certificate
by Result~\ref{res:boolmob}.
\end{result}
\noindent Formalized:
\lref{BooleanMobiusCarry.lean}{734}{exists_normalized_support_fraction_iff_exists_booleanMobiusCarry},
with the certificate packaged as
\lrefx{BooleanMobiusCarry.lean}{589}{BooleanMobiusCarryCertificate} and the
quotient normalized at index zero
(\lrefx{BooleanMobiusCarry.lean}{579}{carryQuotient}). A worked fixture keeps
the coordinates honest: the support $\{2,3\}$ has value
$\tfrac13+\tfrac17=\tfrac{10}{21}$
(\lrefx{BooleanMobiusCarry.lean}{824}{erdosSupportSeries_support23_eq_ten_div_twenty_one}),
its orbit is the pure period-six cycle $10,20,19,17,13,26$ with multiplier
$21$ (\lrefx{BooleanMobiusCarry.lean}{863}{carryOrbit23_isTempered}), and the
M\"obius transform of the carry quotient recovers exactly $\{2,3\}$
(\lrefx{BooleanMobiusCarry.lean}{900}{mobius_carryOrbit23_recovers_support}).

These are coordinates, not by themselves a strategy: nothing here excludes
infinite Boolean carry paths, and no finite search is promoted to a
completeness argument.

\subsection{Residue wraps, reciprocal mass, and unbounded carries}

Suppose now the support series equals $p/(2^c v)$ with $v$ odd and $A$
containing a positive element. Rationality makes the scaled tails integers:
$u(N)=v\,T_{f_A}(c+N)$ is a positive integer, satisfies the carry recurrence
$u(N+1)+v\,f_A(c+N+1)=2\,u(N)$, and reduces modulo $v$ to the doubling
residues $r_N=p\,2^N\bmod v$
(\lrefx{RationalSupportCarrySkeleton.lean}{736}{exists_shifted_odd_tail_nat_states_of_support_fraction}).
Doubling a least residue either stays under $v$ or wraps past it exactly once,
so each step emits a wrap digit $k_N=\lfloor2r_N/v\rfloor\in\{0,1\}$, and over
any cycle length $h$ with $2^h\equiv1\pmod v$ the repetend identity holds:
\[
\sum_{N<h}r_N \;=\; v\cdot w,
\qquad
w=\sum_{N<h}k_N,
\]
with $w\ge1$ when $p$ is coprime to $v>1$. Formalized:
\lref{RationalSupportCarrySkeleton.lean}{209}{sum_doublingResidue_eq_mul_wrapCount}
and \lrefx{RationalSupportCarrySkeleton.lean}{219}{doublingWrapCount_pos};
residues, digits, and counts are
\lrefx{RationalSupportCarrySkeleton.lean}{73}{doublingResidue},
\lrefx{RationalSupportCarrySkeleton.lean}{77}{doublingWrapDigit},
\lrefx{RationalSupportCarrySkeleton.lean}{81}{doublingWrapCount}.

\begin{result}[one-wrap classification]\label{res:onewrap}
Let $v>1$ be odd, $p$ coprime to $v$, and $h\ge1$ with $2^h\equiv1\pmod v$.
If the cycle of $p$ wraps exactly once, then $v=2^h-1$ and the starting
residue is a power of two: $r_0=2^a$ for some $a<h$.
\end{result}
\noindent The one-wrap cycles are thus exactly the Mersenne repetends
$2^a/(2^h-1)$. Formalized:
\lref{RationalSupportCarrySkeleton.lean}{454}{one_doublingWrap_classification},
from the generic closed-cycle form
\lrefx{RationalSupportCarrySkeleton.lean}{392}{one_wrap_cycle_classification},
and instantiated at the true order $h=\operatorname{ord}_v(2)$
(\lrefx{RationalSupportCarrySkeleton.lean}{518}{one_oddOrder_doublingWrap_classification},
with \lrefx{RationalSupportCarrySkeleton.lean}{475}{oddDoublingOrder}). The
classification is algebraic; a finite check of $446$ coprime starting residues
across twelve modulus/order rows is retained as kernel-checked validation, not
as the proof
(\lrefx{RationalSupportCarrySkeleton.lean}{2327}{orderWrapValidatedStartCount_eq_446},
\lrefx{RationalSupportCarrySkeleton.lean}{2331}{orderWrapValidation446_passes},
\lrefx{RationalSupportCarrySkeleton.lean}{2337}{orderWrap_minima_table_passes}).

The analytic bridge is Ces\`aro averaging. Write $\rho(A)=\sum_{a\in A}1/a$
for the reciprocal mass
(\lrefx{RationalSupportCarrySkeleton.lean}{858}{reciprocalMass}). When the
reciprocal family is summable, the mean of $f_A(1),\dots,f_A(N)$ converges to
$\rho(A)$ --- the density of the multiples of $a$, summed over the support ---
and the mean of the scaled tails has the same limit
(\lref{RationalSupportCarrySkeleton.lean}{989}{tendsto_supportCoeff_mean_reciprocalMass},
\lrefx{RationalSupportCarrySkeleton.lean}{1112}{tendsto_binaryCoeffTail_supportCoeff_mean_reciprocalMass}).
The divergent case is carried as an explicit disjunction, not through the
convention that a divergent \texttt{tsum} is zero
(\lrefx{RationalSupportCarrySkeleton.lean}{864}{ReciprocalMassDivergesOrAtLeast}).
Each tail state splits exactly as $u(N)=v\,e_N+r_N$ with integral excess
$e_N=\lfloor u(N)/v\rfloor\ge0$
(\lrefx{RationalSupportCarrySkeleton.lean}{557}{oddTailExcess}).

\begin{result}[order-wrap lower bound, with exact excess mean]\label{res:orderwrap}
Let $A$ and $p/(2^c v)$ be as above with $v>1$ odd, and let $w$ be the wrap
count of $p$ over one cycle of length $h=\operatorname{ord}_v(2)$. Then
$\sum_{a\in A}1/a$ diverges or $\rho(A)\ge w/h$; if moreover $p$ is coprime to
$v$, then $w\ge1$, so $\rho(A)\ge1/\operatorname{ord}_v(2)$. In the summable
case the bound is the periodic part of an exact identity: the Ces\`aro mean of
the excess $e_N$ converges automatically, and
\[
\rho(A)\;=\;\frac{w}{h}\;+\;\lim_{N\to\infty}\frac1N\sum_{n<N}e_n .
\]
\end{result}
\noindent Formalized:
\lref{RationalSupportCarrySkeleton.lean}{1519}{support_fraction_oddOrder_wrapRatio_le_reciprocalMass},
\lref{RationalSupportCarrySkeleton.lean}{1543}{one_div_oddOrder_le_reciprocalMass_of_support_fraction},
the divergence-honest form
\lloc{RationalSupportCarrySkeleton.lean}{1582}\ (\,\idn{support_fraction_one_div_oddOrder_reciprocalMassDivergesOrAtLeast}\,),
and the excess mean
\lref{RationalSupportCarrySkeleton.lean}{1431}{tendsto_oddTailExcess_mean}
(fraction-facing shifted form
\lrefx{RationalSupportCarrySkeleton.lean}{1613}{exists_shifted_oddTailExcess_mean_of_support_fraction},
named-limit form
\lrefx{RationalSupportCarrySkeleton.lean}{1445}{reciprocalMass_eq_wrapRatio_add_oddTailExcessMean}).
A rational value thus ties the odd part of its denominator to the sparsity of
its support: a summable support with small $\rho(A)$ can only display
denominators whose odd part has small multiplicative order of $2$.

\begin{result}[collision strengthening at common multiples]\label{res:collision}
In the summable case, let $F\subseteq A$ be finite and let $L$ be a common
multiple of $F$. Then
\[
\rho(A)\;\ge\;\frac{w}{h}\;+\;\frac{\lceil(|F|-1)/2\rceil}{L}.
\]
In particular, an infinite support whose value is a dyadic rational $p/2^c$
has $\sum_{a\in A}1/a$ divergent or $\rho(A)>1$.
\end{result}
\noindent The mechanism is a collision: at every positive multiple of $L$ at
least $|F|$ support divisors land on the same index, so $f_A\ge|F|$ there,
while the exact carry equation $f_A(n)+e_n=k_{n-1}+2e_{n-1}$ with
$k_{n-1}\le1$ lets one step absorb at most one unit --- so the preceding
excess must spike to at least $\lceil(|F|-1)/2\rceil$, and the spikes recur
with density $1/L$. Formalized:
\lref{RationalSupportCarrySkeleton.lean}{1924}{booleanCollision_wrap_bound_of_common_multiple}
(denominator-shifted form
\lloc{RationalSupportCarrySkeleton.lean}{1863}\ (\,\idn{shiftedBooleanCollision_wrap_bound_of_common_multiple}\,)),
the carry equation
\lrefx{RationalSupportCarrySkeleton.lean}{1734}{supportCoeff_add_oddTailExcess_succ_eq_wrap_add_two_excess},
and the dyadic corollary
\lref{RationalSupportCarrySkeleton.lean}{2074}{dyadic_support_fraction_reciprocalMass_diverges_or_gt_one}
(via \lloc{RationalSupportCarrySkeleton.lean}{2003}\ (\,\idn{one_lt_reciprocalMass_of_dyadic_support_fraction_of_two_pos_mem}\,)).

\begin{result}[carry states are unbounded over infinite supports]\label{res:unbounded}
Let $A$ be infinite with value $p/(2^c v)$, $v\ge1$. Then the positive integer
carry state $u(N)=v\,T_{f_A}(c+N)$ is unbounded: for every $B$ there is an $N$
with $u(N)>B$.
\end{result}
\noindent The proof picks $2B+1$ positive support elements; at a common
multiple $L$ of all of them, the local inequality $1+v\,|F|\le2\,u(L-c-1)$
pushes the state past $B$. Formalized:
\lref{RationalSupportCarrySkeleton.lean}{2251}{exists_unbounded_shifted_odd_tail_nat_state_of_support_fraction},
from
\lrefx{RationalSupportCarrySkeleton.lean}{2194}{shifted_state_unbounded_of_infinite_support}
and the local inequality
\lrefx{RationalSupportCarrySkeleton.lean}{2102}{one_add_mul_card_le_two_mul_shifted_state};
the $\{2,3\}$ fixture revalidates it at sixty-six common multiples
(\lrefx{RationalSupportCarrySkeleton.lean}{2276}{carryOrbit23_common_multiple_bound_validation66}).
So a rational value over an infinite support cannot run on a bounded carry
alphabet. That closes finite-state readings of the carry system; it does not
close the problem. Nothing in Results~\ref{res:onewrap}--\ref{res:unbounded}
prevents an infinite support from satisfying every one of these constraints.
The repetend identities and divisor averages used here are classical. The
coupled reciprocal-mass lower bounds, collision strengthening, and global
unboundedness statement are retained as exact-source-comparison candidates,
not claimed novelties.

\subsection{Greedy geometry of the \#257 value set}

The fourth module studies the set of all values a base-$2$ support series can
take. For $n\ge1$ put $w_n=1/(2^n-1)$, let $T(n)$ be the mass strictly after
exponent $n$, and let
\[
\mathcal A\;=\;\Bigl\{\,\sum_{n\in A}w_n \;:\; A\subseteq\N,\ 0\notin A\Bigr\},
\]
the \emph{Mersenne achievement set}, the exponent $0$ being normalized away as
analytically invisible. Coding by supports of positive exponents is literally
the kernel's support series
(\lrefx{GreedyAchievementSet.lean}{505}{positiveMersenneSupportValue_eq_erdosSupportSeries};
the set is \lrefx{GreedyAchievementSet.lean}{522}{mersenneAchievementSet}). In
this language the base-$2$ \#257 statement reads: the rational members of
$\mathcal A$ are exactly the finite subset sums. Nothing below decides that.
Kova\v{c} and Tao record the strict separation, uniqueness, and Cantor-set
geometry for these Lambert subseries~\cite{kovactao}; the local contribution is
the executable Lean interface, quantitative enclosure, and death-certificate
packaging over that known geometry.

\begin{result}[superincreasing geometry and greedy membership]\label{res:greedy}
For every $n\ge1$ the weight dominates the whole tail after it, $T(n)<w_n$,
with the two-scale gap $w_n-T(n)=\tfrac23\,4^{-n}+O(8^{-n})$ and explicit
valid $O$-constant $3$. Consequently a real $x$ belongs to $\mathcal A$
exactly when $x\ge0$ and every greedy residual of $x$ is at most the remaining
tail; and normalized support coding is injective, so each member of
$\mathcal A$ has a unique support, which the greedy recursion recovers bit by
bit.
\end{result}
\noindent Formalized:
\lref{GreedyAchievementSet.lean}{212}{mersenneTail_lt_weight},
\lref{GreedyAchievementSet.lean}{434}{mersenneGap_isBigOWith} (explicit bound
\lrefx{GreedyAchievementSet.lean}{409}{mersenneGap_asymptotic_bound}),
\lref{GreedyAchievementSet.lean}{738}{mem_mersenneAchievementSet_iff_greedy_survival},
and
\lref{GreedyAchievementSet.lean}{747}{positiveMersenneSupportValue_injective_normalized},
transferred to the kernel series at
\lrefx{GreedyAchievementSet.lean}{756}{erdosSupportSeries_two_injective_normalized}.
The nonnegativity guard in the membership criterion is necessary: without it a
negative target would survive every level while lying outside $\mathcal A$.

\begin{result}[exact rational runs and finite non-membership certificates]\label{res:death}
The greedy recursion executed in exact rational arithmetic agrees step by step
with the real recursion under casting. A finite decidable certificate --- the
exact rational residual at some level is at least an exact rational upper
enclosure of the remaining tail --- soundly proves non-membership in
$\mathcal A$; the level-one instance shows $3/4\notin\mathcal A$. And a
rational number already known to lie in $\mathcal A$ has computable support
bits.
\end{result}
\noindent Formalized:
\lref{GreedyAchievementSet.lean}{575}{cast_greedyMersenneRemainderRat};
\lref{GreedyAchievementSet.lean}{815}{certifiedGreedyMersenneDeath_not_mem},
with the certificate
\lrefx{GreedyAchievementSet.lean}{809}{CertifiedGreedyMersenneDeath} and the
enclosure \lrefx{GreedyAchievementSet.lean}{776}{mersenneTailUpperRat};
\lref{GreedyAchievementSet.lean}{837}{three_fourths_not_mem_mersenneAchievementSet};
and \lref{GreedyAchievementSet.lean}{911}{rational_member_support_bit_iff}.
The certificates are one-sided: survival through any finite depth proves
nothing about membership, and no decidability of membership in $\mathcal A$ is
asserted. The module also declines to formalize the topological or
measure-theoretic geometry of $\mathcal A$; in particular the phrase ``fat
Cantor set'' is not promoted to a Lean-checked theorem.

\paragraph{What this section does not do.} No statement above excludes a
rational value of the totient constant or of any infinite-support series. The
four modules supply coordinates (Results~\ref{res:rigidity},
\ref{res:boolmob}, \ref{res:carrycert}), constraints
(Results~\ref{res:onewrap}--\ref{res:unbounded}), and a decision geometry for
the value set (Results~\ref{res:greedy}, \ref{res:death}). The closing
argument --- excluding every infinite Boolean carry path, or exhibiting an
unbounded certificate supply --- is exactly what Sections~\ref{sec:eb}
and~\ref{sec:249} leave open.

\section{Index of principal declarations}\label{app:index}

Each row links to the declaration at commit \href{\repobase}{\texttt{59f9cc4}}.

\begin{center}
\footnotesize
\begin{tabular}{@{}p{7.2cm}p{7.8cm}@{}}
\toprule
informal statement & linked Lean module\\
\midrule
$\sum 1/(b^n-1)$ irrational, all $b\ge2$ & \lrefx{CertificateKernel.lean}{7999}{irrational_erdosSum_full_support}\ (\texttt{CertificateKernel})\\
Erd\H{o}s--Borwein $E$ irrational & \lrefx{CertificateKernel.lean}{8006}{irrational_erdosBorwein_series}\ (\texttt{CertificateKernel})\\
factorial support & \lrefx{CertificateKernel.lean}{5706}{irrational_erdosSum_factorial_support}\ (\texttt{CertificateKernel})\\
power-of-two support & \lrefx{CertificateKernel.lean}{5730}{irrational_erdosSum_two_pow_support}\ (\texttt{CertificateKernel})\\
pairwise-coprime supports & \lrefx{CertificateKernel.lean}{10447}{irrational_erdosSupportSeries_pairwise_coprime}\ (\texttt{CertificateKernel})\\
$L(\mu)=1/2$ & \lrefx{MersenneLambertLadder.lean}{527}{tsum_moebius_div_two_pow_sub_one_eq_half}\ (\texttt{MersenneLambertLadder})\\
$L(\ph)=2$ & \lrefx{MersenneLambertLadder.lean}{515}{tsum_totient_div_two_pow_sub_one_eq_two}\ (\texttt{MersenneLambertLadder})\\
positive lift $L(A)=S$ & \lrefx{CertificateKernel.lean}{18116}{tsum_primWeight_div_two_pow_sub_one_eq_totient_series}\ (\texttt{CertificateKernel})\\
M\"obius-square lens & \lrefx{CertificateKernel.lean}{18125}{totient_series_eq_half_add_moebius_mersenne_square}\ (\texttt{CertificateKernel})\\
$S=\tfrac12+\Pr(\gcd=1)$ & \lrefx{CertificateKernel.lean}{18228}{totient_series_eq_half_add_visible_coprime_pairs}\ (\texttt{CertificateKernel})\\
$S\ne p/q$ for $q\le\Qzero$ & \lrefx{CertificateKernel.lean}{18055}{tsum_totient_div_pow_two_ne_ratCast_of_den_le_79639646646701375323355774875831053}\ (\texttt{CertificateKernel})\\
tail-period reduction & \lrefx{TotientTailPeriodKiller.lean}{386}{irrational_totient_series_of_certificate_supply}\ (\texttt{TotientTailPeriodKiller})\\
diagonal reduction & \lrefx{LcmDiagonalReduction.lean}{92}{irrational_totient_series_of_lcm_diagonal_certificate_supply}\ (\texttt{LcmDiagonalReduction})\\
cone flatness from rationality & \lrefx{LcmConeFlatness.lean}{90}{rational_totient_series_forces_lcm_cone_flatness}\ (\texttt{LcmConeFlatness})\\
certificate completeness & \lrefx{LcmConeFlatness.lean}{316}{exists_certifiedKill_iff_tail_diff_notMem_int}\ (\texttt{LcmConeFlatness})\\
joint cone refuter & \lrefx{LcmConeNonflat.lean}{126}{exists_nonintegral_pair_of_coneNonflatCert}\ (\texttt{LcmConeNonflat})\\
verified finite instances & \lrefx{TotientTailPeriodKiller.lean}{396}{certifiedKill_all_small}\ (\texttt{TotientTailPeriodKiller})\\
tail-orbit rationality criterion & \lrefx{GenericTailOrbitRigidity.lean}{260}{binaryCoeffSeries_rational_iff_exists_temperedBinaryOrbit}\ (\texttt{GenericTailOrbitRigidity})\\
Boolean carry certificate $\Leftrightarrow$ support fraction & \lrefx{BooleanMobiusCarry.lean}{734}{exists_normalized_support_fraction_iff_exists_booleanMobiusCarry}\ (\texttt{BooleanMobiusCarry})\\
one-wrap classification & \lrefx{RationalSupportCarrySkeleton.lean}{454}{one_doublingWrap_classification}\ (\texttt{RationalSupportCarrySkeleton})\\
order-wrap mass bound & \lrefx{RationalSupportCarrySkeleton.lean}{1543}{one_div_oddOrder_le_reciprocalMass_of_support_fraction}\ (\texttt{RationalSupportCarrySkeleton})\\
unbounded carry states & \lrefx{RationalSupportCarrySkeleton.lean}{2251}{exists_unbounded_shifted_odd_tail_nat_state_of_support_fraction}\ (\texttt{RationalSupportCarrySkeleton})\\
greedy membership criterion & \lrefx{GreedyAchievementSet.lean}{738}{mem_mersenneAchievementSet_iff_greedy_survival}\ (\texttt{GreedyAchievementSet})\\
\bottomrule
\end{tabular}
\end{center}

\clearpage
\begin{thebibliography}{6}
\small
\bibitem{erdos1948} P.~Erd\H{o}s, \emph{On arithmetical properties of Lambert series},
  J. Indian Math. Soc. (N.S.) \textbf{12} (1948), 63--66.
\bibitem{erdos1968} P.~Erd\H{o}s, \emph{On the irrationality of certain series},
  Math. Student \textbf{36} (1968), 222--226.
\bibitem{nesterenko} Yu.~V. Nesterenko, \emph{Modular functions and transcendence questions},
  Mat. Sb. \textbf{187} (1996), no.~9, 65--96; English transl., Sb. Math. \textbf{187} (1996), no.~9, 1319--1348.
\bibitem{lucatachiya} F.~Luca and Y.~Tachiya, \emph{Irrationality of Lambert series associated
  with a periodic sequence}, Int. J. Number Theory \textbf{10} (2014), no.~3, 623--636.
\bibitem{kovactao} V.~Kova\v{c} and T.~Tao, \emph{On several irrationality problems for
  Ahmes series}, arXiv:2406.17593 (2024).
\bibitem{wanggrauribas} H.~Wang and J.~M. Grau Ribas, \emph{Positive dyadic density for
  rational weighted binary expansions}, arXiv:2606.24972 (2026).
\end{thebibliography}

\end{document}
